\chapter{Introduction}
\label{chap:introduction}

\todo{Write full introduction.}

Modern recommender systems achieve state-of-the-art accuracy by learning dense,
entangled latent representations of users and items.
While these black-box representations optimise prediction metrics, they are
fundamentally difficult for users to understand or correct, and for developers
to audit.
This lack of transparency is particularly critical in the domain of
\emph{Points of Interest} (POI) recommendation, where user preferences are
strongly tied to geographic and temporal context and change dynamically.

\section{Motivation}

\todo{Expand motivation.}

The goal of this thesis is to apply \emph{Sparse Autoencoders} (SAE) to the
latent space of a collaborative filtering model trained on Yelp data.
By enforcing sparsity, the model is encouraged to learn monosemantic,
human-interpretable features that can simultaneously explain recommendations
and serve as interactive ``knobs'' for user-driven steering.

\section{Research questions}

\begin{enumerate}
  \item Can SAE-derived sparse features maintain recommendation quality comparable
        to dense baselines?
  \item Do the learned features correspond to semantically interpretable
        preference dimensions (e.g.\ cuisine type, ambience, price range)?
  \item Can users effectively steer recommendations by manipulating these features
        through an interactive interface?
\end{enumerate}

\section{Thesis outline}

\Cref{chap:related_work} reviews related work on collaborative filtering,
sparse representations, and interpretable recommendation.
\Cref{chap:methodology} presents the proposed SAE-CF architecture.
\Cref{chap:experiments} describes the experimental setup and results.
\Cref{chap:conclusion} concludes and outlines future work.
